\chapter{高斯态与协方差抽样}
\label{app:gaussian}

本附录给出可直接实现的高斯态检查与抽样流程。采用
\begin{equation}
  \bm R=(Q_1,P_1,\ldots,Q_M,P_M)^T,
  \qquad
  \comm{\hat R_j}{\hat R_k}=\ii\Omega_{jk},
\end{equation}
其中
$\Omega=\bigoplus_{n=1}^M\begin{pmatrix}0&1\\-1&0\end{pmatrix}$。

\section{均值、协方差和物理性}

高斯态由
\begin{equation}
  \bar{\bm R}=\qavg{\hat{\bm R}},
  \qquad
  V_{jk}=\frac12\qavg{
  \Delta\hat R_j\Delta\hat R_k+
  \Delta\hat R_k\Delta\hat R_j}
\end{equation}
完全确定。必要且充分的物理性条件是
\begin{equation}
  V+\frac{\ii}{2}\Omega\succeq0.
  \label{eq:appendix-gaussian-physicality}
\end{equation}
真空的 $V=\identity/2$ 恰好饱和该不等式。

Williamson 定理保证存在辛矩阵 $S$ 使
\begin{equation}
  V=S\diag(\nu_1,\nu_1,\ldots,\nu_M,\nu_M)S^T,
\end{equation}
其中辛本征值 $\nu_j\ge1/2$。精确算术下，$\ii\Omega V$ 的本征值为实数对 $\{+\nu_j,-\nu_j\}$；这里的“取绝对值”是对这些本征值逐个取模，绝不是求矩阵 $\ii\Omega V$ 的奇异值。后者一般不等于辛本征值。

更稳定的数值做法是先由对称特征分解构造正定平方根 $V^{1/2}$，再形成 Hermitian 矩阵
\begin{equation}
  K=\ii V^{1/2}\Omega V^{1/2}=K^\dagger.
  \label{eq:hermitian-symplectic-spectrum}
\end{equation}
$K$ 的本征值同样按 $\pm\nu_j$ 成对出现，可用 Hermitian 本征求解器得到。实际实现应检查正负配对误差，并以正本征值作为 $\nu_j$。普通矩阵特征值全为正并不足以保证式\eqref{eq:appendix-gaussian-physicality}；对 $V>0$，该条件等价于所有 $\nu_j\ge1/2$。

\section{常见单模协方差}

\begin{center}
\begin{tabular}{lll}
\toprule
状态 & 均值 & 协方差 \\
\midrule
真空 & $(0,0)$ & $\diag(1/2,1/2)$ \\
相干态 $\alpha_0$ & $(\sqrt2\Re\alpha_0,\sqrt2\Im\alpha_0)$ & $\identity_2/2$ \\
热态 $\bar n$ & $(0,0)$ & $(\bar n+1/2)\identity_2$ \\
实参数压缩真空 & $(0,0)$ & $\tfrac12\diag(\ee^{-2r},\ee^{2r})$ \\
\bottomrule
\end{tabular}
\end{center}
本书固定 $r>0$ 压缩 $Q$ 轴。最后一行对应
\begin{equation}
  \alpha=\cosh r\,\beta-\sinh r\,\beta^*,
\end{equation}
于是 $Q_\alpha=\ee^{-r}Q_\beta$、$P_\alpha=\ee^rP_\beta$，与第4章完全一致。若使用相反号的压缩算符约定，必须同时交换压缩轴或令 $r\to-r$，不能只改抽样公式的一处符号。

\section{Cholesky 抽样算法}

给定 $\bar{\bm R},V$：
\begin{enumerate}
  \item 以双精度读入并对称化 $V\leftarrow(V+V^T)/2$；
  \item 检查维数为偶数、单位与正交变量顺序一致；
  \item 按\cref{eq:hermitian-symplectic-spectrum} 计算辛本征值，检查 $\pm$ 配对，并确认最小正值在容差内不小于 $1/2$；
  \item 作 $V=LL^T$ 的 Cholesky 分解；
  \item 为每条轨迹生成 $2M$ 个独立标准正态数 $\bm\xi$；
  \item 令 $\bm R=\bar{\bm R}+L\bm\xi$，再转换为
  $\alpha_n=(Q_n+\ii P_n)/\sqrt2$；
  \item 用独立统计代码回测均值、协方差和目标正规排序量。
\end{enumerate}
合法有限能量高斯态的 $V$ 是正定的。纯态条件是所有 $\nu_j=1/2$，等价于 $\det V=4^{-M}$，并不意味着普通特征值接近零；真空纯态 $V=I/2$ 的条件数就是 $1$。数值病态通常来自强压缩或变量缩放不良，例如单模压缩态的普通特征值为 $\ee^{\mp2r}/2$。此时可用对称本征分解构造平方根；任何特征值裁剪都要记录阈值，并重新检查物理性、协方差和目标正规排序矩。

\section{复协方差形式}

对零均值复模式向量 $\bm\alpha$，定义
\begin{equation}
  C_{ij}=\ensavg{\alpha_i\alpha_j^*},
  \qquad
  M_{ij}=\ensavg{\alpha_i\alpha_j}.
\end{equation}
$C$ 是复协方差的正常块，$M$ 是反常块；两者都来自 Wigner 样本的对称排序矩，不能把“正常块”误读成已经完成正规排序修正。它们与实协方差等价。例如
\begin{align}
  \ensavg{Q_iQ_j}&=\Re(C_{ij}+M_{ij}),\\
  \ensavg{P_iP_j}&=\Re(C_{ij}-M_{ij}),\\
  \frac12\ensavg{Q_iP_j+P_jQ_i}&=\Im(M_{ij}-C_{ij}),
\end{align}
具体下标转置取决于向量排序，实施时应通过单模真空和已知压缩态核对符号。

对于无压缩独立热模式，$C_{ij}=(\bar n_i+1/2)\delta_{ij}$、$M=0$。突然淬火后 $M$ 一般非零；若代码只保存每模占据 $C_{ii}-1/2$，就丢失了成对关联。

\section{高斯重构何时是精确的}

只有当目标量子态为高斯态时，$(\bar{\bm R},V)$ 才确定全部状态信息。对非高斯态，用同一均值和协方差生成高斯样本只是二阶矩闭合：它精确匹配所有一次和二次对称排序矩，但一般改变四阶累积量、粒子数尾部、Wigner 负值与相干干涉。非线性哈密顿量会把这些未匹配矩反馈到低阶观测量，因此“初始协方差正确”不是任意时间正确的保证。

采用这种重构时，应把验收分为两层：第一层检查样本确实重现所构造的 $\bar{\bm R},V$；第二层用精确小系统、不同的矩匹配方案或目标高阶矩，量化“把原态替换成高斯态”所产生的状态误差。第一层通过不能替代第二层。

\section{BdG 模式抽样清单}

本书采用与原子场展开负号一致的约定
\begin{equation}
  \bm w_\nu=\begin{pmatrix}\bm u_\nu\\-\bm v_\nu\end{pmatrix},
  \qquad
  \bar{\bm w}_\nu=\tau_x\bm w_\nu^*
  =\begin{pmatrix}-\bm v_\nu^*\\\bm u_\nu^*\end{pmatrix},
  \qquad
  S=\begin{pmatrix}U&-V^*\\-V&U^*\end{pmatrix}.
  \label{eq:appendix-bdg-sign-and-complete-basis}
\end{equation}
这里 $U=(\bm u_1,\ldots,\bm u_M)$、$V=(\bm v_1,\ldots,\bm v_M)$；
$\bm w_\nu$ 是正能、正辛范数列，$\bar{\bm w}_\nu$ 是其负能、负辛范数
粒子--空穴伙伴；两组列共同组成完整矩阵 $S$。对每个 BdG 基：
\begin{enumerate}
  \item 使用 $\Sigma_z=\diag(\identity,-\identity)$ 检查
  $\bm w_\nu^\dagger\Sigma_z\bm w_\mu=\delta_{\nu\mu}$；
  \item 检查正能列与式\eqref{eq:appendix-bdg-sign-and-complete-basis}的伙伴成对；
  \item 检查 $S^\dagger\Sigma_zS=\Sigma_z$ 以及
  $S^{-1}=\Sigma_zS^\dagger\Sigma_z$；
  \item 在简并子空间内作辛正交化，而不是依赖求解器任意返回的线性组合；
  \item 排除或单独处理 Goldstone 零模；
  \item 由 $b_\nu$ 的热 Wigner 方差生成式\eqref{eq:old-bdg-direct-field-sampling}；
  \item 从原子场样本反投影回 $b_\nu$，检查协方差闭环；
  \item 比较重构的投影核与目标 $\delta_{\mathcal C}$。
\end{enumerate}

\section{数守恒近似的验收}

若采用固定总粒子数抽样，至少报告
\begin{equation}
  \ensavg{N},\quad \Var(N),\quad
  \ensavg{N_0},\quad \ensavg{N_{\rm dep}},\quad
  \Cov(N_0,N_{\rm dep}).
\end{equation}
还应比较重缩放前后的一体密度矩阵与异常协方差。固定每条轨迹的原始
对称范数 $\mathcal N_{\mathrm{sym}}$ 只保证对应的 $N_W$ 方差为零；它不能
保证其他矩仍对应目标量子态。

\section{抽样误差的预期尺度}

对独立高斯样本，协方差元素的标准误按 $N_s^{-1/2}$ 缩小，但系数取决于四阶矩。零均值单变量的样本方差近似满足
\begin{equation}
  {\rm SD}(s^2)\simeq V\sqrt{\frac{2}{N_s-1}}.
\end{equation}
因此协方差验收阈值应随目标方差和轨迹数变化，不能对所有模式使用同一个绝对误差。低占据模式宜检查标准化残差，高占据模式同时检查相对误差。
